\section{نتیجه‌گیری}

این مثال نشان می‌دهد که چگونه می‌توان با استفاده از روش‌های یادگیری عمیق و شبکه عصبی کانولوشنی، عمر باقی‌مانده باتری‌های لیتیوم-یون را پیش‌بینی کرد. در این فرایند، داده‌های مربوط به ۴۰ باتری پردازش و به قالبی مناسب برای ورود به شبکه عصبی تبدیل شدند. سپس داده‌ها به سه مجموعه آموزش، اعتبارسنجی و آزمون تقسیم شدند و برای آموزش و ارزیابی مدل مورد استفاده قرار گرفتند. در نهایت، مدل آموزش‌دیده توانست عمر باقی‌مانده باتری‌ها را با میانگین خطای نسبی حدود ۱۸ درصد تخمین بزند.
